\section{The definite identity and the periodized arithmetic kernel}

% BEGIN I75K1
\paragraph{I75K1. Finiteness and positivity for the marked definite mass.}
For the level $M>0$, the original self-dual definite lattices form
finitely many classes under the rational unitary group $D_H(M)$.
The standard lattice is among them. Each stabilizer is finite and
nonempty, so each reciprocal stabilizer weight is positive.

Indeed for each lattice $L$ take the integral positive definite
symmetric form $b=\operatorname{Tr}_{E/\mathbb Q}h$ on its
underlying free abelian group. Its trace dual is
$\mathfrak D_E^{-1}L$ by Hermitian self-duality. Local freeness gives
$[L^{\#}:L]=|\operatorname{disc}E|^n$. Thus the rank and determinant
of the integral trace lattice are fixed. These data give finitely
many isometry classes. For example, nonzero vectors have length
at least one and the covolume is fixed, so Minkowski's second theorem
bounds every successive minimum. Take a full independent set of
bounded length. Its integral Gram matrix has finitely many
possibilities, and the index of its span is bounded by the
covolume ratio. There are finitely many overlattices of that span
of bounded index, since one common factorial clears all of them.

On each resulting trace lattice, multiplication by any chosen
integral basis element $a$ of $\mathcal O_E$ is an integral
endomorphism of operator norm at most $\max_\sigma|\sigma(a)|$,
by the orthogonal Minkowski decomposition of $b$. There are
finitely many such endomorphisms in a fixed basis. An isometry
preserving these operators is $\mathcal O_E$-linear and preserves
$h$, since $b(ax,y)=\operatorname{Tr}(a\,h(x,y))$ determines
$h(x,y)$ by nondegeneracy of the trace pairing. It extends to a
rational unitary isometry. We have therefore obtained finiteness
of the rational unitary isometry classes in the prescribed set.

At primes above $M$ the lattice is the standard local lattice by
the agreement condition. Every isometry between two of the
prescribed lattices belongs to the group integral with integral
inverse at these primes. Reduction of this group modulo $M$ has
finite image. Its kernel is $D_H(M)$: entrywise local divisibility
of $g-1$ by $M$ is equivalent, by clearing the remaining
denominators, to the source's integer multiplier prime to $M$.
Consequently passage to this kernel splits each class into only
finitely many classes. It contains the identity, is closed under
products and inverse, and acts on the same set of lattices.
The stabilizer injects into the isometry group of a positive
integral trace lattice, a discrete subgroup of a compact group,
finite also by the bounded possible images of a basis. The standard
lattice has all the defining properties; the zero-rank case gives
one lattice and the identity automorphism. This verifies the mass
properties used by the Hecke and cusp arguments.
% END I75K1

% BEGIN I75K2
\paragraph{I75K2. The definite identity for every admitted family.}
Let $E$ be a CM field with CM type $\Psi$, let $n$ be even, and
let $4d<n$. At degree $d$ the weighted definite theta series at
zero elliptic arguments equals the weighted Eisenstein series there
for every input of the original $\mathrm{SWAt}$ assertion.
G75W1--G75W3 construct the Hermitian Weil action used here, with its
original theta transformation, inner product and Heisenberg covariance.
For the archimedean period reduction we use A75B2.

Induct on $d$ with $E,\Psi,n$ fixed. Degree zero is Z75D1. For
$d>0$ the strict range holds also in degree $d-1$. Fix positive
$M,R$ and one $(M,R)$-admitted family $\phi$. By I75K1 and the
support, periodicity and equivariance requirements it has finite
supremum norm over all definite lattices in $V(M)$ and all test
arguments. B75S1--B75S4 apply the induction hypothesis to the
restricted and moved family at every rational boundary direction.
C75E4 supplies absolute convergence, including a bound on a whole
high cone, and W75A2 retains admission at a common level under
the period reduction. Thus A75B3 bounds the invariant norm
for all families at any one fixed level, uniformly in their norm
and in the period.

The same bound must be available for all twists of $\phi$ before
the Hecke contraction can be used. Choose the single integer $c$
of Z75D6 for degree $d$. Apply A75B3 at $(M,Rc)$, independently
of the twisting lattice $X$. Z75D6 transforms the resulting bound
back to $I(X)^{n/2}J(Z)^{n/2}|D_X(Z)|$, with a constant times
the supremum norm of $\phi$. It holds for every allowed $X,Z$.
P75H3 supplies the Hecke identity for these same twists. Hence the
finite supremum is at most itself times the factor of Z75D5,
strictly smaller than one. It is zero. Taking $X=\mathcal O_E^d$
gives $D_\phi(Z)=0$ for every $Z$. All choices in this induction
come after $M,R,\phi$ have been fixed, except for the constants
which were explicitly required to be uniform. This proves the
zero-argument identity at the specified nonnegative degrees.
% END I75K2

% BEGIN I75K3
\paragraph{I75K3. The differentiated identity and spanning on a supplied system.}
Return to a source level system $(\Delta,\mathcal L,\Phi,\mu)$
of positive integer level $M_0$. Use the original valid period
frames, finite Weil and doubling transformations, doubled theta
product and graph-coordinate law of O74E2--O74E4. Fix even
$n>4r$, a pair of lattices $L,L'\in\mathcal L$, and a finite
invariant family $(P_j)$ of admissible row-highest harmonic
polynomials of the prescribed weight. With $C(L,L')$ the stabilizer-weighted class
sum of $\langle\Phi(L,\Lambda),\Phi(L',\Lambda)\rangle$, the
arithmetic identity in the spanning argument is
\[
\begin{split}
 a\sum_{\gamma\in\Delta} C(L,\gamma L')
       \overline{\mu(\gamma,L')} K_P(x,\gamma y)
    =\sum_{[\Lambda],j}\operatorname{wt}(M_0,\Lambda)\,
      \Theta_{\Lambda j}(L,x)
      \overline{\Theta_{\Lambda j}(L',y)} .
\end{split}
\]
The scalar $a\ne0$ is the one from O74E2. It was chosen before the
lattices, frames, tests in the system and harmonic family were fixed.

Indeed the admission clause of this system supplies one $R$ for the
doubled tests for the pair $L,L'$. Apply I75K2 to this very family
at $(M_0,R)$ in degree $r$. It gives $H(Z,0,0)=0$ for every
allowed period, where $H$ is the weighted theta--Eisenstein
difference of O74E4. O74E5 and C75E4 give its heat equations and
holomorphy. Since $r\le n$, H75L1--H75L2 kill its differentiated
Gaussian pullback at the doubled period for $x,y$. Therefore the
weighted theta and Eisenstein sides have equal derivatives.
On the Eisenstein side O74E1--O74E4 remove the singular terms,
count each graph once and justify differentiation of its series.
On the theta side the doubled product formula turns the derivative
into exactly the displayed finite sum; the conjugated polynomial
acts on the independent second variable. This proves (ICL.6) on
the stated inputs with the actual class weights and phases.

For even $n\ge 4[E:\mathbb Q]r+4$ in the weight range of ICL14,
with the specified nonzero invariant harmonic family, K74A4 supplies (ICL.7) with a
strictly positive real scalar. For the original lattice tests F74N3
gives strict diagonal dominance. Applying the inverse class-sum
operator of ICL12 to a bounded holomorphic form, the reproducing
integral and the absolutely convergent unfolding of ICL13--ICL14
now express that form as a finite linear combination of
$\Theta_{\Lambda j}$. One set of coefficients works at all lattices
and frames in the supplied system. To obtain the geometric
short-section span of ICL16 one still needs the actual principal
charts and level-system data, and the componentwise identification
of these analytic theta functions with the algebraic sections in
(ICL.12). The original doubled test, product and graph laws remain
explicit inputs in the arithmetic identity above.
% END I75K3
